To classify news articles into one of four categories (world, sports, business, and sci/tech), we use the AG News dataset from Hugging Face \cite{zhang2015character}. This dataset consists of news articles with associated labels, titles, and descriptions. Each class contains 30,000 training samples and 1,900 testing samples, resulting in a total of 120,000 training samples and 7,600 testing samples.

We used the official Hugging Face split (120,000 train and 7,600 test). The training set was split into a new training set (of size 108,000) and a development set (of size 12,000) using stratified sampling with a fixed random seed of 67 to ensure reproducibility. The final split sizes are: train 108,000, dev 12,000, and test 7,600.

\paragraph{Preprocessing.} We combined the title and description of each news article into a single text input. We cleaned the text by HTML unescaping to convert HTML entities back to their original characters (like \texttt{\#39;} to \texttt{'}) and removed any remaining HTML tags. We also unescaped \LaTeX entities (\verb|\\| etc.), collapsed backslashes and normalized whitespace. We converted the labels from 1-based to 0-based indexing (i.e., world: 0, sports: 1, business: 2, sci/tech: 3) for easier handling in the classification models. We combine the title and description into a single text input in the format: \texttt{Title: <title> Description: <description>} to provide a clear structure for the models to learn from.

\paragraph{Tokenization.} We used BERT Tokenizer, which is based on WordPiece tokenization, to tokenize the text inputs. Similar to BPE, WordPiece is a subword tokenization method that can effectively handle out-of-vocabulary words (OOV) and reduce the vocabulary size while still capturing meaningful subword units. We used the pre-trained BERT tokenizer (\textit{distilbert-base-uncased} from Hugging Face, which has a vocabulary size of 30,522 tokens. Both BERT and LSTM models were trained using the same tokenizer to ensure consistency in input representation.